\problemname{Tågmani}
\noindent
Det är väldigt välkänt att Julia gillar serien \textit{På spåret}. Då kanske det inte kommer som en
överraskning att hon gillar att åka tåg. Tyvärr har hon nyligen spårat ur lite, och denna
tidigare oskyliga hobby har förvandlats till ett maniskt behov att åka tåg. Problemet är att
det är svårt att åka tåg så mycket som hon skulle vilja, eftersom resorna tar lång tid.
Även om hon inte kan åka med alla tåg, så försöker hon resa med så många som möjligt.
Efter att hon suttit timtals i sitt rum och försökt skapa den bästa reseplanen bestämer
du dig för att hjälpa henne. Även om mängden tåg hon åker inte är hälsosamt, så föredrar
du att hon får åka tåg istället för att planera sin tågåkning.

Julia bor i staden Lönnköping. Lönnköping har $N$ stationer numrerade från $1$ till $N$.
För att resa mellan dessa finns det två transportmedel: elscootrar och tåg. Det finns
så många elscootrar i staden att hon kan påbörja en resa på elscooter från vilken station som helst.
Hon kan bara åka elscooter längs med de $N-1$ vägarna i städarna. Man kan färdas från
vilken station till vilken annan som helst hjälp av vägarna dessa, alltså utgör vägnätverket ett
\textit{träd}. Dessa vägar har olika längd, och det tar därför olika lång tid att färdas längs med
vägarna om hon använder en scooter. \\
Hon kan även färdas med de $K$ nästkommande tågavgångarna. Varje tåg har ett strikt schema:
det avgår från en viss station vid en viss tid, och anländer vid dess slutstation efter ett antal sekunder.
Om Julia befinner sig vid en station samma sekund som ett tåg avgår kan hon välja att åka med det och sedan
stiga av vid slutstationen (hon kan inte förmå sig att stiga av ett tåg innan hon verkligen måste).
Om ett tag anländer vid en station som sekund som ett annat avgår hinner hon göra bytet om hon så vill.

Givet hur vägnätverket ser ut och de $K$ nästa tågavgångarna, hur många tåg kan Julia åka
med om hon optimalt planerar sin resa? Eftersom planen är en bit in i framtiden är hon inte
säker på var hon ska börja, så hon vill att du beräknar största antalet tåg hon kan åka
med för varje möjlig startstation. Du kan anta att hon alltid start vid tid $0$.


\section*{Indata}
Den första raden i indatan innehåller två heltal, $N$ och $K$ ($1 \le N, K \le 10^5$), antalet hållplatser
det finns i Lönnköping och antalet tågavgångar som du ska betrakta.

De följande $N-1$ beskriver vägnätverket. Varje av dessa rader innehåller heltalen $u,v,s$
($1 \leq u, v \leq N$, $1 \leq s \leq 10^9$), som betyder att det finns en väg mellan
stationerna med nummer $u$ och $v$. Det tar $s$ sekunder att åka
från $u$ till $v$ eller $v$ till $u$ med vägen om man använder en elscooter.

Därefter följer $K$ rader som beskriver tågavgångarna. Varje av dessa rader innehåller heltalen
$u,v,t,d$ ($1 \leq u,v \leq N$, $u \neq v$, $1 \leq t, d \leq 10^9$), vilket betyder att
vid $t$ så avgår det ett tåg från station $u$ till station $v$, som anländer vid station $v$
vid tid $t+d$. Det är garanterat att alla $t$ är unika. 

\section*{Utdata}
För varje möjlig startstation $1 \leq i \leq N$, skriv ut antalet tåg hon kan åka med
om hon börjar sin resa vid den stationen. Skriv ut dessa separerade med mellanslag i
ökande ordningen av $i$, det vill säga först $i=1$, sen $i=2$ osv.

\section*{Poängsättning}
Din lösning kommer att testas på en mängd testfallsgrupper.
För att få poäng för en grupp så måste du klara alla testfall i gruppen.

\text{Graden} av en station är antalet vägar som den är sammankopplad med.
\noindent
\begin{tabular}{| l | l | p{12cm} |}
  \hline
  \textbf{Grupp} & \textbf{Poäng} & \textbf{Gränser} \\ \hline
  $1$    & $10$       & $N, T, K \leq 50$ \\ \hline
  $2$    & $5$        & $N, K \leq 50$ \\ \hline
  $3$    & $15$       & $N \leq 1000, K \leq 200$ \\ \hline
  $4$    & $15$       & $N, K \leq 2000$ \\ \hline
  $5$    & $10$       & $K \leq 2000$ och det minsta $t$:et är större än summan av alla $s$. \\ \hline
  $6$    & $20$       & $K \leq 30000$ och ingen station har grad större än 3. \\ \hline
  $7$    & $25$       & Inga ytterligare begränsningar. \\ \hline
\end{tabular}

%\section*{Exempelfall}

